% Generated draft: requires mathematical and source review.
\documentclass[11pt]{article}
\usepackage[margin=25mm]{geometry}
\usepackage{fontspec}
\setmainfont{Latin Modern Roman}
\usepackage{amsmath,amssymb,mathtools,unicode-math}
\setmathfont{Latin Modern Math}
\usepackage{enumitem,needspace,xurl,hyperref}
\setlength{\parindent}{0pt}
\setlength{\parskip}{5pt}
\setlength{\emergencystretch}{3em}
\newcommand{\N}{\mathbb{N}}
\newcommand{\Z}{\mathbb{Z}}
\newcommand{\Q}{\mathbb{Q}}
\newcommand{\R}{\mathbb{R}}
\newcommand{\C}{\mathbb{C}}
\newcommand{\F}{\mathbb{F}}
\newcommand{\A}{\mathbb{A}}
\newcommand{\PP}{\mathbb{P}}
\newcommand{\E}{\mathbb{E}}
\newcommand{\eps}{\varepsilon}
\newcommand{\dd}{\,\mathrm{d}}
\begin{document}
\section*{UnsolvedMath ID: 30002493}
\textbf{Source problem code:} catalogue entry
\subsection{Definitions and mathematical statement}
Let \(B\) denote the basic Newton stratum in the setting intended by the source: a Hodge-type Shimura variety or a moduli space of shtukas. Let \(\Gamma\) be the intended indexing set of Mirković–Vilonen cycles, and let \(\{Z_\gamma\}_{\gamma\in\Gamma}\) be the associated family of algebraic cycles. The supplied record does not specify the geometric data defining \(B\), the particular Mirković–Vilonen cycles used, or the construction and ambient space of the cycles \(Z_\gamma\).

For the intended cohomological degree and twist, write \(H\) for the cohomology space in which the cycle classes lie, \(K\) for its coefficient field, and \(\operatorname{cl}(Z_\gamma)\in H\) for the corresponding cycle classes. Write \(T_{\mathrm{pred}}\subseteq H\) for the subspace of Tate classes predicted by the intended Langlands conjecture. These are placeholders for data not specified in the record; in particular, the arithmetic Galois action defining Tate classes and the precise Langlands prediction remain unspecified.

The source asks whether, for its intended geometric and cohomological data,
\[
\operatorname{span}_{K}\{\operatorname{cl}(Z_\gamma):\gamma\in\Gamma\}
= T_{\mathrm{pred}}.
\]
Here the span consists of finite \(K\)-linear combinations. If the family contains cycles in different codimensions, this question must be interpreted in the corresponding separate cohomological degrees and twists, or in a specified graded space; the record does not choose between these formulations. The displayed equality records the intended spanning question, rather than an established assertion. A fully specified conjecture requires the missing geometric data, cycle construction, cohomology theory, and definition of the predicted target.
\subsection{Short English statement}
For the basic Newton stratum of a Hodge-type Shimura variety or a moduli space of shtukas, the source describes cycles indexed by Mirković–Vilonen cycles. It asks whether their cohomology classes span all Tate classes predicted by the intended Langlands conjecture. The precise geometric setting, cycle classes, and predicted subspace require clarification.
\subsection{Sources}
\begin{itemize}
\item Local UnsolvedMath registry, supplied complete source record: permanent ID \(30002493\), title \emph{Algebraic Cycles Spanning Predicted Tate Classes}, exact source problem code \texttt{catalogue entry}, in the Oberwolfach Reports Open Problems collection. The registry gives proposed year 2014 and supplies no named proposer.
\item Supplied report citation: \emph{Langlands Correspondence and Constructive Galois Theory} (2014), Oberwolfach Reports, pp. 287--334, DOI \texttt{10.4171/owr/2014/05}. Source URL and literature-source label ``OWR report'': \url{https://doi.org/10.4171/owr/2014/05}. The supplied citation identifies the report as a whole and gives no problem-specific page.
\item Supplied literature-source label ``cycles on Shimura varieties via geometric Satake'': \url{https://arxiv.org/abs/1707.05700}.
\item Supplied literature-source label ``Tate cycles on unitary Shimura varieties'': \url{https://doi.org/10.2140/ant.2017.11.2213}.
\item Supplied literature assessment, checked \texttt{2026-08-21}: the registry labels the problem partially solved. Its assessment says that the reported Mirković–Vilonen-cycle families and their geometric partition are established, and that linear independence realizes the expected Tate classes in generic function-field type-A cases. It says that general spanning of all predicted Tate classes in the broader Shimura-variety setting remains unresolved. This is supplied dated curation, not independently verified here; the record explicitly requests independent verification of time-sensitive status and cited claims.
\item The registry labels the statement \texttt{corrected\_verified} and reports that it was checked against the official Oberwolfach report, with source hypotheses and notation restored. This verification claim is supplied registry metadata, not an independent source check performed for this draft.
\end{itemize}
\end{document}
